\clearpage
\section{Error Analysis}\label{sec:errors}
\subsection{Sampling and annotation design}
The analysis targets errors of the final DistilBERT + LoRA model. Its 1,859
misclassified test rows reduce to 1,840 unique review texts after deduplication.
Using seed 42, we sample 100 unique reviews in two strata:
40 from 702 reviews wrong for every model (\code{all5}), and 60 from the other
1,138 DistilBERT errors. Sampling fractions are 0.0570 and 0.0527. The labels
therefore describe the deep model's errors, rather than the complete error
distribution of any classical pipeline.

One annotator assigns one or more categories, placing the main reason first:
sarcasm, negation, mixed opinions, \code{long_truncated}, possible
\code{label_noise}, or other. Two annotation files, A and B, are halves labeled
by the same person; they do not represent independent annotators. Counts are
calculated from the final 100-row annotation table.

\subsection{Category counts and stratified estimates}
Table~\ref{tab:categories} and Figure~\ref{fig:categories} summarize the labels.
Main categories sum to 100; any-category counts overlap. Sixty-two reviews have
multiple categories. Mixed opinions are the most common main explanation
(47/100) and appear in 72/100 reviews. Negation occurs in 24 but is the main
reason in only five, illustrating its frequent interaction with other phenomena.

\begin{table}[htbp]\centering
\caption{Categories of 100 sampled DistilBERT errors. ``Any category'' counts a review even when the category is secondary.}\label{tab:categories}
\begin{tabular}{lrr}\toprule
Category & Main category & Any category\\\midrule
\tableinput{tables/error_categories.tex}
\bottomrule\end{tabular}
\end{table}
\reportfigure{person4_rId27.png}{0.98}{Main and overlapping category counts for the sampled errors. These are sample counts, not model-wide error rates.}{fig:categories}

To account for the sampling design, the estimated share of category $c$ in
the 1,840 unique DistilBERT errors is
\begin{equation}
 \widehat p_c=\sum_{h\in\{\mathrm{all5},\mathrm{other}\}}
 \frac{N_h}{1840}\frac{n_{h,c}}{n_h}.
\end{equation}
Here $N_h$ and $n_h$ are the stratum population and sample sizes. We use
a stratified variance estimate with finite-population corrections
to construct approximate 95\% intervals. For mixed opinions as the main
category, the weighted estimate is 47.14\%, with interval 37.55--56.73\%.
The estimate is close to the unweighted 47\%, but the interval remains wide.
It concerns unique deep-model errors, not all test reviews.

Possible label noise is the main category for 10 of 40 \code{all5} reviews and
7 of 60 other reviews. Sarcasm is the main category for 1/40 and 8/60,
respectively. These small subgroup counts are descriptive. The two annotation
halves contain 28 and 19 mixed-main reviews, but the difference cannot be
attributed to different annotators.

Only approximately 33\% of TF-IDF + LR errors, 34\% of TF-IDF + LinearSVC
errors and 28\% of either Word2Vec model's errors are also DistilBERT errors.
Thus even the overlap with the sampled population covers only part of each
classical model's mistakes. Applying the sample's category percentages to all
classical errors would be unjustified.

\subsection{Representative examples}
Table~\ref{tab:error_examples} gives excerpts from our sampled error cases.
Every listed full review is a DistilBERT error; its prediction is the opposite
of the binary gold label. Excerpts are explanatory evidence, not independently
labeled benchmark examples.

\begin{table}[htbp]\centering\small
\caption{Selected qualitative examples. Gold labels refer to complete reviews; ``label noise'' is an annotator hypothesis.}\label{tab:error_examples}
\begin{tabularx}{\linewidth}{@{}l l Y@{}}\toprule
Review / category & Gold & Excerpt and interpretation\\\midrule
\code{test_02554} / sarcasm & Negative &
``hilarious antics ensue (I wish)''; positive wording is reversed by the aside.\\
\code{test_14247} / negation & Positive &
``it's not too bad''; negation softens a negative adjective.\\
\code{test_00383} / negation & Negative &
``It's not half as good any of those films.''; comparative negation expresses criticism.\\
\code{test_13076} / mixed & Positive &
``really stupid and lame for the first half''\ldots ``good for the last half''; the opinion changes over the review.\\
\code{test_02576} / mixed & Negative &
``Most of this film was okay''\ldots ``you should go watch something else.''\\
\code{test_14218} / truncation & Positive &
``Undoubtedly it has been a major success'' occurs in the unseen middle.\\
\code{test_06593} / label noise & Negative &
``one of the strongest Russian movie for the past few years''\ldots ``8/10'' suggests tension with the dataset label.\\
\code{test_24629} / other & Positive &
``PLEASE PICK UP THE SHOW!!!!'' expresses support through anger at cancellation.\\\bottomrule
\end{tabularx}
\end{table}

The examples illustrate scope and target problems: negative language can refer
to a character or a cancellation rather than the film's quality. Conversely,
a locally positive phrase may be subordinate to an overall negative verdict.
The sarcasm fragment in the demo is predicted positive by every model, whereas
the negation fragment is correctly interpreted by DistilBERT. Fragment behavior
does not contradict errors on the complete source reviews.

\subsection{Truncation and possible label noise}
Seventeen of the 100 sampled errors exceed the 510-wordpiece review budget,
but only two are annotated \code{long_truncated}. In \code{test_14218}, strong
praise is removed from the middle while the visible end is largely critical.
In \code{test_22750}, ``I liked Brain of Blood a lot'' is removed, but the
retained end still includes ``I liked it''. These observations make truncation
a plausible explanation for one case and a possible contributor to the second;
length alone does not establish causation. No counterfactual full-text model
was evaluated to prove that either prediction would change.

Twenty sampled reviews receive a possible label-noise category, 17 as the main
reason. IMDB labels originate in ratings rather than a separate manual sentiment
annotation of the review text. The apparent ``8/10'' contradiction in
\code{test_06593} is suggestive, but the exports have no independent rating
field with which to verify it. Gold labels are retained for every reported
metric; suspected disagreement is not used to improve scores retrospectively.

\subsection{Heuristic checks and AI assistance}
We also compare keyword heuristics with the final
any-category labels. A contrast cue has the largest reported Cohen's kappa,
0.3353, for mixed opinions. A broad negation cue fires in 93/100 reviews but
has precision 0.2581 for negation, close to the 0.2400 prevalence. The truncation
flag finds both annotated truncation cases, yet its precision is only 2/17,
or 0.1176. These checks do not justify turning heuristic flags into category
error rates. In particular, a length flag is a measurable input property,
not an explanation of an incorrect prediction.

Annotation used AI suggestions. Suggestions were shown for 97 reviews; 69
were accepted unchanged and 28 edited. Eighteen edits only removed an unsure
flag. Three reviews were initially labeled before suggestions existed, and
one of these, \code{test_04366}, was revised after its suggestion was seen.
Consequently, 98 final labels were set with the suggestion available. Main
categories agree with the AI in 97/100 reviews, with kappa 0.9566. This is
AI-versus-final agreement, not inter-annotator reliability. Nine final records
remain marked unsure. The strong agreement means that the category distribution
may partly reflect the suggestion system's tendencies; it is not a blind,
independent human replication.

\subsection{Threats to validity and next steps}
The study has several limits that affect interpretation:
\begin{itemize}
\item Classical configuration changes involved observing test scores, whereas
      deep-model selection used validation only. Future comparisons should freeze
      every pipeline before one test pass.
\item DistilBERT receives eight tuning configurations and external pretraining;
      classical models receive limited tuning. Equal-budget searches and
      multiple seeds are needed to separate these effects.
\item Internal test duplicates remain; near-duplicates are not excluded. A
      duplicate-aware sensitivity analysis would assess their influence.
\item Classical filtering and mean pooling can remove compositional information;
      the transformer still truncates long inputs. Preprocessing ablations and
      a full-document or window-based comparison would test these explanations.
\item The qualitative sample covers deep-model errors only and has one
      AI-assisted annotator. A second independent annotator, explicit adjudication
      of unsure cases and samples from each model would improve reliability.
\item All data are English movie reviews. Transfer to other review domains,
      languages and genuinely neutral opinions is unmeasured.
\end{itemize}
\FloatBarrier
